\documentclass[12pt,a4paper]{article}

% ── Packages ──────────────────────────────────────────────────────────────────
\usepackage[a4paper, margin=2.5cm]{geometry}
\usepackage[bahasa]{babel}
\usepackage[utf8]{inputenc}
\usepackage[T1]{fontenc}
\usepackage{graphicx}
\usepackage{booktabs}
\usepackage{array}
\usepackage{tabularx}
\usepackage{enumitem}
\usepackage{listings}
\usepackage{xcolor}
\usepackage{amsmath}
\usepackage{amssymb}
\usepackage{parskip}
\usepackage{titlesec}
\usepackage{fancyhdr}
\usepackage{caption}
\usepackage{float}
\usepackage{hyperref}
\usepackage[space]{grffile}  % handle spasi/titik dalam nama file gambar

% ── Hyperref setup ────────────────────────────────────────────────────────────
\hypersetup{
    colorlinks=true,
    linkcolor=black,
    urlcolor=blue,
    citecolor=black
}

% ── Listing style for Rust code ───────────────────────────────────────────────
\definecolor{rustbg}{RGB}{248,248,248}
\definecolor{rustkw}{RGB}{175,0,219}
\definecolor{rustcmt}{RGB}{106,153,85}
\definecolor{ruststr}{RGB}{206,145,120}
\definecolor{rustnum}{RGB}{181,206,168}

\lstdefinelanguage{Rust}{
    keywords={fn, let, mut, struct, impl, enum, match, if, else, for, loop,
               while, return, use, pub, mod, self, Self, true, false, as,
               in, ref, move, where, type, const, static, extern, crate,
               super, trait, dyn, async, await, Box, Vec, String, Option,
               Some, None, Result, Ok, Err, Arc, Mutex},
    keywordstyle=\color{rustkw}\bfseries,
    comment=[l]{//},
    commentstyle=\color{rustcmt}\itshape,
    stringstyle=\color{ruststr},
    morestring=[b]",
    sensitive=true
}

\lstset{
    language=Rust,
    backgroundcolor=\color{rustbg},
    basicstyle=\ttfamily\small,
    breaklines=true,
    frame=single,
    rulecolor=\color{gray!40},
    numbers=left,
    numberstyle=\tiny\color{gray},
    numbersep=8pt,
    tabsize=4,
    showstringspaces=false,
    captionpos=b
}

% ── Section formatting ─────────────────────────────────────────────────────────
\titleformat{\section}{\large\bfseries}{\thesection.}{0.5em}{}
\titleformat{\subsection}{\normalsize\bfseries}{\thesubsection.}{0.5em}{}
\titleformat{\subsubsection}{\normalsize\bfseries\itshape}{\thesubsubsection.}{0.5em}{}
\titlespacing{\section}{0pt}{2ex plus .5ex minus .2ex}{1ex}
\titlespacing{\subsection}{0pt}{1.5ex}{0.5ex}
\titlespacing{\subsubsection}{0pt}{1ex}{0.3ex}

% ── Header / Footer ───────────────────────────────────────────────────────────
\pagestyle{fancy}
\fancyhf{}
\rhead{\small ETS Algoritma Pemrograman -- ELISA}
\lhead{\small Kelompok 4}
\cfoot{\thepage}

% Halaman TOC tanpa header fancy
\fancypagestyle{plain}{%
    \fancyhf{}
    \cfoot{\thepage}
    \renewcommand{\headrulewidth}{0pt}
}

% ── Document ──────────────────────────────────────────────────────────────────
\begin{document}

% ── Cover Page ────────────────────────────────────────────────────────────────
\begin{titlepage}
    \centering
    \vspace*{2cm}

    {\LARGE\bfseries Evaluasi Tengah Semester (ETS)\\[0.3em]
    Algoritma Pemrograman\par}

    \vspace{1cm}
    {\large\bfseries Pengembangan Sistem Instrumentasi Pengukuran dan Kontrol\\
    Berbasis Rust\par}

    \vspace{0.8cm}
    {\large\bfseries
    ``ELISA Electrical Monitoring and Intelligent Safety Alarm\\
    Sistem Monitoring Energi Listrik Rumah / UMKM Berbasis Rust''\par}

    \vfill
    \includegraphics[width=5cm]{LAMBANG ITS.png}
    \vfill

    {\normalsize
    \begin{tabular}{ll}
        \multicolumn{2}{c}{\bfseries Kelompok 4:} \\[8pt]
        Devan Ale Satria       & 2042251047 \\[3pt]
        Sheilka Janice Leandro & 2042251086 \\
    \end{tabular}\par}

    \vspace{2cm}
\end{titlepage}

\newpage
\thispagestyle{plain}
\tableofcontents
\newpage

% ══════════════════════════════════════════════════════════════════════════════
\section{Pendahuluan}
% ══════════════════════════════════════════════════════════════════════════════

\subsection{Latar Belakang}

Keselamatan instalasi listrik di rumah tangga dan usaha kecil menengah (UMKM)
merupakan isu yang sering diabaikan. Berdasarkan data Kementerian ESDM, sebagian
besar kebakaran di Indonesia disebabkan oleh korsleting listrik akibat kebocoran
arus yang tidak terdeteksi secara dini. Kebocoran arus listrik selain dapat memicu
kebakaran juga membahayakan keselamatan jiwa; arus sebesar 15\,mA sudah cukup
menyebabkan kontraksi otot pada manusia, dan di atas 50\,mA dapat menyebabkan
fibrilasi jantung.

Sistem monitoring arus bocor yang ada di pasaran umumnya berbentuk alat fisik
sederhana tanpa kemampuan pemantauan secara \emph{real-time} dan tidak memberikan
informasi analitik kepada pengguna. Oleh karena itu, pada project ETS ini
dikembangkan \textbf{ELISA} (\textit{Electrical Monitoring and Intelligent Safety
Alarm}), sebuah sistem monitoring energi listrik berbasis perangkat lunak Rust
yang terintegrasi dengan \emph{dashboard} web, mampu mendeteksi kebocoran arus
secara \emph{real-time}, dan memberikan peringatan bertingkat sesuai standar
PUIL 2011.

\subsection{Tujuan}

\begin{itemize}
    \item Merancang dan mengimplementasikan sistem monitoring energi listrik berbasis bahasa pemrograman Rust.
    \item Menerapkan konsep pemrograman berbasis objek (OOP) menggunakan \texttt{struct} dan \texttt{impl}.
    \item Mengimplementasikan komputasi numerik berupa \emph{Moving Average}, standar deviasi, dan akumulasi energi.
    \item Menyajikan hasil monitoring secara \emph{real-time} melalui \emph{dashboard} web yang dapat diakses dari \emph{browser}.
\end{itemize}

\subsection{Manfaat}

Sistem ini dapat digunakan sebagai prototipe perangkat lunak untuk sistem
monitoring kelistrikan di rumah tangga maupun UMKM, sebelum diintegrasikan ke
\emph{hardware} nyata seperti ESP32 dengan sensor SCT-013.

% ══════════════════════════════════════════════════════════════════════════════
\section{Studi Kasus Instrumentasi}
% ══════════════════════════════════════════════════════════════════════════════

\subsection{Topik}

Sistem Monitoring Energi Listrik deteksi kebocoran arus pada instalasi listrik
rumah/UMKM dengan tiga cabang MCB (\textbf{DAPUR}, \textbf{KAMAR},
\textbf{R.TAMU}).

\subsection{Sensor yang Digunakan}

\textbf{SCT-013} (\textit{Split Core Current Transformer / CT Clamp})

SCT-013 adalah sensor arus non-invasif yang bekerja berdasarkan prinsip
transformator arus (\textit{current transformer}). Sensor ini dipasang menjepit
kabel tanpa perlu memutus instalasi. \textit{Output} berupa tegangan AC yang
proporsional terhadap arus yang mengalir, kemudian dibaca oleh ADC
mikrokontroler ESP32.

Dalam sistem ELISA, setiap cabang MCB menggunakan dua buah SCT-013: satu
dipasang di kabel fasa dan satu lagi di kabel netral. Prinsip deteksi kebocoran
mengikuti hukum Kirchhoff:
\[
    I_{\text{bocor}} = I_{\text{fasa}} - I_{\text{netral}}
\]

Dalam kondisi normal, arus yang masuk melalui fasa sama dengan arus yang keluar
melalui netral, sehingga selisihnya mendekati nol. Jika ada kebocoran arus ke
bumi (\textit{ground}), arus netral akan lebih kecil dari arus fasa, dan
selisihnya merepresentasikan besar arus bocor.

\subsection{Threshold Berdasarkan PUIL 2011 dan Referensi Fisiologis}

\begin{table}[H]
    \centering
    \caption{Threshold status kebocoran arus berdasarkan PUIL 2011}
    \begin{tabular}{>{\bfseries}lcccc}
        \toprule
        Kondisi  & Rentang $I_{\text{bocor}}$ & LED    & Buzzer \\
        \midrule
        AMAN     & $< 15\,\text{mA}$          & Hijau  & OFF    \\
        WASPADA  & $15\text{--}22\,\text{mA}$ & Kuning & OFF    \\
        BAHAYA   & $> 22\,\text{mA}$          & Merah  & ON     \\
        \bottomrule
    \end{tabular}
\end{table}

Nilai 15\,mA dipilih karena merupakan ambang batas di mana seseorang masih bisa
melepaskan konduktor secara mandiri. Nilai 22\,mA dipilih sebagai batas atas
WASPADA sebelum memasuki zona bahaya fibrilasi jantung.

\subsection{Parameter yang Dimonitor}

Setiap cabang MCB memiliki parameter monitoring berikut:

\begin{itemize}
    \item Tegangan ($V$) dari ADC ESP32
    \item Arus fasa $I_{\text{fasa}}$ (mA) dari SCT-013 pertama
    \item Arus netral $I_{\text{netral}}$ (mA) dari SCT-013 kedua
    \item Arus bocor $I_{\text{bocor}} = I_{\text{fasa}} - I_{\text{netral}}$ (mA), hasil kalkulasi
    \item Daya aktif $P = V \times I_{\text{fasa}}$ (W)
    \item \emph{Moving Average} arus fasa MA-5 (mA)
    \item Standar deviasi arus fasa $\sigma$
    \item Akumulasi energi total (kWh)
\end{itemize}

% ══════════════════════════════════════════════════════════════════════════════
\section{Algoritma dan Flowchart}
% ══════════════════════════════════════════════════════════════════════════════

\subsection{Algoritma Sistem}

\subsubsection*{Algoritma Utama}

\begin{enumerate}
    \item \textbf{MULAI}
    \item Inisialisasi sistem (3 sensor: DAPUR, KAMAR, R.TAMU)
    \item Jalankan web server di port 7878 (thread terpisah)
    \item Tampilkan menu: \texttt{[1] Input Manual \quad [2] Simulasi \quad [3] Injeksi Bocor \quad [4] Keluar}
    \item Baca pilihan pengguna
    \begin{enumerate}[label=\alph*.]
        \item \textbf{JIKA} pilihan = 1:\\
              Untuk tiap sensor: baca tegangan, $I_{\text{fasa}}$, netral dari terminal\\
              Validasi data $\to$ hitung $I_{\text{bocor}}$ $\to$ update aktuator $\to$ perbarui \emph{dashboard}
        \item \textbf{JIKA} pilihan = 2:\\
              Generate data simulasi otomatis\\
              Tampilkan di \emph{dashboard}, tanya lanjut atau tidak\\
              Jika tidak lanjut: reset sensor ke standby AMAN
        \item \textbf{JIKA} pilihan = 3:\\
              Cek \emph{baseline} (jika belum ada data, isi dulu semua cabang ke AMAN)\\
              Tampilkan status tiap cabang saat ini\\
              Pilih skenario: WASPADA / BAHAYA / Reset cabang / Reset semua\\
              Hanya ubah cabang yang dipilih, cabang lain tetap
        \item \textbf{JIKA} pilihan = 4:\\
              Keluar dari program
    \end{enumerate}
    \item \textbf{ULANGI TERUS}
    \item \textbf{SELESAI}
\end{enumerate}

\subsubsection*{Algoritma Kalkulasi Per Sensor}

\begin{verbatim}
FUNGSI update_all(sensor):
  validate(sensor)      -> cek tegangan & arus dalam batas wajar
  compute_leak(sensor)  -> I_bocor = I_fasa - I_netral
  push_history(sensor)  -> tambahkan I_fasa ke riwayat MA-5
  update_actuator(sensor)-> set LED & buzzer sesuai status
  process(controller, status)
                        -> catat jumlah kejadian waspada/bahaya
  update_energy(system) -> E += (P_total / 1000) x 1 jam
\end{verbatim}

\subsubsection*{Algoritma Deteksi Status}

\begin{verbatim}
FUNGSI status(sensor):
  JIKA is_valid = FALSE  -> kembalikan BAHAYA
  JIKA leak_ma >= 22.0   -> kembalikan BAHAYA
  JIKA leak_ma >= 15.0   -> kembalikan WASPADA
  LAINNYA                -> kembalikan AMAN
\end{verbatim}

\subsection{Flowchart}

\begin{figure}[H]
    \centering
    \includegraphics[width=0.5\textwidth, height=0.88\textheight, keepaspectratio]{Flowchart alprog.drawio.png}
    \caption{Flowchart sistem ELISA}
    \label{fig:flowchart}
\end{figure}

% ══════════════════════════════════════════════════════════════════════════════
\section{Penjelasan Program Rust}
% ══════════════════════════════════════════════════════════════════════════════

\subsection{Struktur Program}

\begin{table}[H]
    \centering
    \caption{Struktur modul program ELISA}
    \small
    \begin{tabularx}{\textwidth}{>{\bfseries}lX}
        \toprule
        Bagian & Isi / Fungsi \\
        \midrule
        Konstanta Sistem       & Menyimpan parameter sistem seperti threshold kebocoran, tegangan nominal, dan batas arus maksimum. \\[4pt]
        LeakStatus (enum)      & Representasi level status kebocoran listrik: Aman, Waspada, dan Bahaya. \\[4pt]
        Actuator (struct)      & Representasi aktuator seperti LED indikator dan buzzer alarm. \\[4pt]
        Sensor (struct + impl) & Menyimpan data sensor tiap cabang MCB beserta logika pembacaan dan analisis kondisi kebocoran. \\[4pt]
        Controller (struct + impl) & Mengatur pencatatan kejadian, pengambilan keputusan, dan pesan tindakan sistem. \\[4pt]
        MonitoringSystem (struct + impl) & Mengoordinasikan seluruh sensor dan controller agar sistem monitoring berjalan terintegrasi. \\[4pt]
        render\_dashboard      & Fungsi untuk menghasilkan tampilan dashboard berbasis HTML pada web server. \\[4pt]
        Helper Input           & Kumpulan fungsi pendukung untuk membaca dan memvalidasi input dari terminal. \\[4pt]
        input\_manual          & Mode untuk memasukkan data sensor secara manual melalui terminal. \\[4pt]
        input\_simulasi        & Mode simulasi otomatis menggunakan data yang dihasilkan program. \\[4pt]
        inject\_kebocoran      & Mode demonstrasi untuk mensimulasikan kondisi kebocoran listrik pada sistem. \\[4pt]
        main                   & Entry point program yang menjalankan web server dan loop menu utama sistem ELISA. \\
        \bottomrule
    \end{tabularx}
\end{table}

\subsection{Fitur Utama Program}

\textbf{Input Data Sensor}\\
Program mendukung tiga mode input: manual (pengguna mengetik nilai), simulasi
otomatis (data deterministik acak berbasis \emph{seed}), dan injeksi kebocoran
(demo skenario waspada/bahaya).

\textbf{Percabangan}\\
Digunakan secara ekstensif pada penentuan status (\texttt{LeakStatus::from\_ma}),
validasi sensor (\texttt{validate}), dan seluruh logika menu.

\textbf{Perulangan}\\
\texttt{loop} pada menu utama, \texttt{for} pada iterasi sensor, dan \texttt{loop}
pada \emph{helper} input terminal untuk validasi input pengguna.

\textbf{Validasi Data}\\
Fungsi \texttt{validate()} memeriksa bahwa tegangan berada di rentang 100--260\,V
dan arus tidak melebihi batas maksimum sensor. Jika tidak valid, sensor otomatis
berstatus BAHAYA.

\textbf{Dashboard Web Real-Time}\\
Web server \texttt{tiny\_http} berjalan di \emph{thread} terpisah dan menyajikan
\emph{dashboard} HTML yang \emph{auto-refresh} setiap 3 detik. Dashboard
menampilkan kartu per sensor dengan indikator LED, grafik riwayat arus, dan tabel
alarm.

\textbf{Concurrency dengan \texttt{Arc} dan \texttt{Mutex}}\\
Data sistem dibungkus \texttt{Arc<Mutex<MonitoringSystem>{}>} untuk berbagi
\emph{state} secara aman antara \emph{thread} utama (input terminal) dan
\emph{thread} web server.

\subsection{Potongan Kode Penting}

\textbf{Kalkulasi arus bocor:}
\begin{lstlisting}[language=Rust, caption={Fungsi compute\_leak}]
fn compute_leak(&mut self) {
    self.leak_ma = (self.current_fasa_ma
        - self.current_netral_ma).max(0.0);
}
\end{lstlisting}

\textbf{Validasi sensor:}
\begin{lstlisting}[language=Rust, caption={Fungsi validate}]
fn validate(&mut self) {
    self.is_valid = self.voltage >= 100.0
        && self.voltage <= 260.0
        && self.current_fasa_ma >= 0.0
        && self.current_fasa_ma <= MAX_CURRENT_MA
        && self.current_netral_ma >= 0.0
        && self.current_netral_ma <= MAX_CURRENT_MA;
}
\end{lstlisting}

\textbf{Penentuan status keseluruhan sistem:}
\begin{lstlisting}[language=Rust, caption={Fungsi overall\_status}]
fn overall_status(&self) -> LeakStatus {
    if self.sensors.iter().any(|s| s.status() == LeakStatus::Bahaya) {
        LeakStatus::Bahaya
    } else if self.sensors.iter().any(|s| s.status() == LeakStatus::Waspada) {
        LeakStatus::Waspada
    } else {
        LeakStatus::Aman
    }
}
\end{lstlisting}

% ══════════════════════════════════════════════════════════════════════════════
\section{Konsep OOP}
% ══════════════════════════════════════════════════════════════════════════════

\subsection{Penerapan Struct dan Impl}

Program ELISA menggunakan pendekatan OOP Rust melalui \texttt{struct} sebagai
representasi objek dan \texttt{impl} sebagai definisi \emph{method}-nya. Berikut
tiga objek utama yang dipersyaratkan:

\textbf{Objek Sensor}\\
Sensor merepresentasikan satu cabang MCB dengan semua atributnya. \emph{Method}
yang dimiliki antara lain \texttt{validate()}, \texttt{compute\_leak()},
\texttt{push\_history()}, \texttt{moving\_average()}, \texttt{std\_deviation()},
\texttt{status()}, dan \texttt{update\_actuator()}.

\begin{lstlisting}[language=Rust, caption={Struct Sensor}]
struct Sensor {
    name: String,
    voltage: f32,
    current_fasa_ma: f32,
    current_netral_ma: f32,
    leak_ma: f32,
    history_fasa: VecDeque<f32>,
    is_valid: bool,
    actuator: Actuator,
}
\end{lstlisting}

\textbf{Objek Controller}\\
Controller bertanggung jawab atas logika keputusan dan pencatatan \emph{event}
per cabang. Method utama: \texttt{process()} (mencatat jumlah waspada/bahaya)
dan \texttt{action\_message()} (menghasilkan pesan tindakan sesuai status).

\begin{lstlisting}[language=Rust, caption={Struct Controller}]
struct Controller {
    name: String,
    total_warnings: u32,
    total_dangers: u32,
}
\end{lstlisting}

\textbf{Objek MonitoringSystem}\\
MonitoringSystem adalah sistem induk yang mengkoordinasikan seluruh sensor dan
controller. Method utama: \texttt{update\_all()}, \texttt{total\_power()},
\texttt{dominant\_branch()}, \texttt{update\_energy()}, dan
\texttt{overall\_status()}.

\begin{lstlisting}[language=Rust, caption={Struct MonitoringSystem}]
struct MonitoringSystem {
    sensors: Vec<Sensor>,
    controllers: Vec<Controller>,
    total_energy_kwh: f32,
    reading_count: u32,
}
\end{lstlisting}

\textbf{Objek Aktuator (tambahan)}\\
Actuator merepresentasikan \emph{output} fisik \emph{hardware} (LED dan buzzer)
yang tertanam di dalam tiap Sensor. Ini mendemonstrasikan konsep komposisi objek
-- Sensor \emph{memiliki} Actuator sebagai atribut.

\subsection{Enkapsulasi}

Setiap \texttt{struct} hanya mengekspos \emph{method} yang diperlukan. Logika
internal seperti perhitungan \texttt{leak\_ma} dan \texttt{is\_valid} dikelola
sepenuhnya oleh \emph{method}-method milik Sensor, tanpa perlu diakses langsung
dari luar.

\subsection{Polimorfisme melalui Enum}

\texttt{LeakStatus} (enum dengan tiga varian: Aman, Waspada, Bahaya) berperan
sebagai mekanisme polimorfisme. Method \texttt{label()} dan \texttt{color()} pada
enum menghasilkan \emph{output} berbeda tergantung varian yang aktif, tanpa
memerlukan kondisi \texttt{if-else} berulang di luar \texttt{struct}.

\begin{lstlisting}[language=Rust, caption={Method color pada LeakStatus}]
fn color(&self) -> &str {
    match self {
        Self::Aman     => "#16a34a",
        Self::Waspada  => "#d97706",
        Self::Bahaya   => "#dc2626",
    }
}
\end{lstlisting}

% ══════════════════════════════════════════════════════════════════════════════
\section{Implementasi Komputasi Numerik}
% ══════════════════════════════════════════════════════════════════════════════

Program ELISA mengimplementasikan empat metode komputasi numerik yang semuanya
relevan dengan pengolahan sinyal sensor listrik.

\subsection{Kalkulasi Arus Bocor (Diferensial Arus)}

\textbf{Prinsip:} Berdasarkan hukum Kirchhoff tentang arus -- dalam rangkaian
tertutup normal, arus masuk (fasa) harus sama dengan arus keluar (netral). Jika
ada kebocoran ke \emph{ground}, selisihnya adalah arus bocor.

\textbf{Formula:}
\[
    I_{\text{bocor}} = \max(I_{\text{fasa}} - I_{\text{netral}},\; 0)
\]

\textbf{Implementasi:}
\begin{lstlisting}[language=Rust, caption={Fungsi compute\_leak}]
fn compute_leak(&mut self) {
    self.leak_ma = (self.current_fasa_ma
        - self.current_netral_ma).max(0.0);
}
\end{lstlisting}

Fungsi \texttt{.max(0.0)} memastikan hasil tidak negatif -- \emph{noise} ADC
bisa menyebabkan $I_{\text{netral}}$ sedikit lebih besar dari $I_{\text{fasa}}$
dalam kondisi ideal, yang secara fisik berarti tidak ada kebocoran.

\subsection{Moving Average (MA-5)}

\textbf{Prinsip:} \emph{Moving average} digunakan untuk meredam \emph{noise}
dari pembacaan ADC. Dengan mengambil rata-rata dari 5 pembacaan terakhir,
fluktuasi sesaat tidak langsung mempengaruhi nilai yang ditampilkan.

\textbf{Formula:}
\[
    \text{MA}_n = \frac{x_n + x_{n-1} + x_{n-2} + x_{n-3} + x_{n-4}}{5}
\]

\textbf{Implementasi:}
\begin{lstlisting}[language=Rust, caption={Fungsi moving\_average}]
fn moving_average(&self) -> f32 {
    if self.history_fasa.is_empty() { return 0.0; }
    self.history_fasa.iter().sum::<f32>()
        / self.history_fasa.len() as f32
}
\end{lstlisting}

\texttt{VecDeque} digunakan sebagai buffer geser -- saat buffer penuh (5 data),
data terlama dihapus otomatis sebelum data baru masuk (\texttt{pop\_front()} +
\texttt{push\_back()}).

\subsection{Standar Deviasi (Sample, n-1)}

\textbf{Prinsip:} Standar deviasi dari riwayat arus fasa digunakan sebagai
indikator kestabilan arus. Nilai $\sigma$ yang besar menandakan arus tidak
stabil, yang dapat menjadi \emph{early warning} meski arus bocor masih di bawah
\emph{threshold}.

\textbf{Formula:}
\[
    \sigma = \sqrt{\frac{\sum_{i=1}^{n}(x_i - \bar{x})^2}{n-1}}
\]

\textbf{Implementasi:}
\begin{lstlisting}[language=Rust, caption={Fungsi std\_deviation}]
fn std_deviation(&self) -> f32 {
    if self.history_fasa.len() < 2 { return 0.0; }
    let avg = self.moving_average();
    let var: f32 = self.history_fasa.iter()
        .map(|x| (x - avg).powi(2))
        .sum::<f32>()
        / (self.history_fasa.len() as f32 - 1.0);
    var.sqrt()
}
\end{lstlisting}

Pembagi $n-1$ (bukan $n$) digunakan karena ini adalah standar deviasi sampel --
data yang tersedia hanya sebagian kecil dari keseluruhan pembacaan sensor
sepanjang waktu.

\subsection{Akumulasi Energi (Numerik Integrasi Euler)}

\textbf{Prinsip:} Energi adalah integral daya terhadap waktu. Karena pembacaan
bersifat diskrit, digunakan metode integrasi numerik Euler \emph{forward} --
mengasumsikan daya konstan selama satu interval pembacaan.

\textbf{Formula:}
\[
    E_{\text{total}} \mathrel{+}= \frac{P_{\text{total}}}{1000} \times \Delta t
    \quad (\Delta t = 1\,\text{jam per pembacaan})
\]

\textbf{Implementasi:}
\begin{lstlisting}[language=Rust, caption={Fungsi update\_energy}]
fn update_energy(&mut self) {
    self.total_energy_kwh +=
        (self.total_power() / 1000.0) * 1.0;
    self.reading_count += 1;
}
\end{lstlisting}

Pembagian 1000 mengkonversi dari Watt ke kilowatt, sehingga satuan hasil adalah
kWh.

% ══════════════════════════════════════════════════════════════════════════════
\section{Hasil Pengujian}
% ══════════════════════════════════════════════════════════════════════════════

\subsection{Pengujian Mode Input Manual}

\textbf{Skenario:} Input data tiga cabang dengan variasi kondisi.

\textbf{Hasil:} Dashboard web menampilkan kartu DAPUR dengan LED hijau, KAMAR
dengan LED kuning, dan R.TAMU dengan LED merah dan buzzer aktif. Tabel alarm di
bawah kartu menampilkan dua entri (KAMAR dan R.TAMU) dengan pesan tindakan yang
sesuai.

\subsection{Pengujian Mode Simulasi Otomatis}

\textbf{Skenario:} Jalankan simulasi 5 iterasi berturut-turut.

Simulasi menggunakan data deterministik berbasis \emph{seed} untuk mereproduksi
kondisi yang bervariasi: 50\% pembacaan menghasilkan status AMAN, 25\% WASPADA,
dan 25\% BAHAYA per cabang. Setiap iterasi memperbarui \emph{dashboard} dan
mengakumulasi data riwayat untuk \emph{Moving Average}.

\textbf{Hasil:} Moving Average MA-5 terlihat mulus dan tidak melonjak drastis
meski ada satu pembacaan ekstrem. Standar deviasi naik ketika ada transisi dari
AMAN ke BAHAYA, mengkonfirmasi fungsinya sebagai \emph{early warning} kestabilan.

\subsection{Pengujian Mode Injeksi Kebocoran}

\textbf{Skenario 1 -- Injeksi tunggal (BAHAYA pada KAMAR):}\\
Pilih menu 3 $\to$ pilih BAHAYA $\to$ pilih KAMAR

\begin{table}[H]
    \centering
    \caption{Hasil pengujian skenario injeksi tunggal}
    \begin{tabular}{lccccc}
        \toprule
        \textbf{Cabang} & \textbf{Tegangan (V)} & \textbf{$I_\text{fasa}$ (mA)}
            & \textbf{$I_\text{netral}$ (mA)} & \textbf{$I_\text{bocor}$ (mA)}
            & \textbf{Status} \\
        \midrule
        DAPUR   & 220 & 800 & 799 & 1.0  & AMAN     \\
        KAMAR   & 220 & 500 & 482 & 18.0 & WASPADA  \\
        R. TAMU & 220 & 350 & 320 & 30.0 & BAHAYA   \\
        \bottomrule
    \end{tabular}
\end{table}

\textbf{Skenario 2 -- Injeksi ganda (WASPADA DAPUR + BAHAYA KAMAR):}\\
Injeksi WASPADA ke DAPUR $\to$ $I_{\text{bocor}}$ DAPUR $\approx$ 18\,mA.
Kembali menu 3, injeksi BAHAYA ke KAMAR $\to$ $I_{\text{bocor}}$ KAMAR
$\approx$ 30\,mA. R.TAMU tetap AMAN.\\
\emph{Dashboard:} DAPUR kuning, KAMAR merah, R.TAMU hijau \checkmark

\textbf{Skenario 3 -- Reset parsial:}\\
Reset cabang DAPUR saja $\to$ DAPUR kembali hijau; KAMAR tetap merah
\checkmark

\textbf{Skenario 4 -- Setelah simulasi, masuk injeksi:}\\
Jalankan simulasi 3x, ketik \texttt{n} untuk reset. Masuk menu 3 $\to$
\emph{baseline} otomatis terisi AMAN. Injeksi BAHAYA ke R.TAMU $\to$ hanya
R.TAMU merah \checkmark

\subsection{Pengujian Validasi Sensor Error}

\textbf{Skenario:} Input tegangan 0\,V (sensor terputus).

\textbf{Hasil:} \texttt{validate()} mengembalikan \texttt{is\_valid = false}.
Status langsung BAHAYA tanpa menunggu perhitungan $I_{\text{bocor}}$. Dashboard
menampilkan label ``Sensor ERROR'' berwarna merah pada cabang tersebut.

% ══════════════════════════════════════════════════════════════════════════════
\section{Analisis Hasil}
% ══════════════════════════════════════════════════════════════════════════════

\subsection{Ketepatan Deteksi Kebocoran}

Sistem berhasil mendeteksi dan mengklasifikasikan kebocoran arus sesuai threshold
PUIL 2011 dengan akurasi 100\% pada semua skenario pengujian. Prinsip diferensial
$I_{\text{fasa}} - I_{\text{netral}}$ terbukti efektif sebagai representasi
perangkat lunak dari cara kerja sensor SCT-013 pada \emph{hardware} nyata.

\subsection{Efektivitas Moving Average}

Moving Average MA-5 terbukti efektif meredam fluktuasi data. Pada pengujian
simulasi, terdapat kasus di mana satu pembacaan menghasilkan arus fasa 850\,mA
disusul pembacaan 750\,mA; tanpa MA, selisih ini terlihat besar di
\emph{dashboard}. Dengan MA-5, nilai yang ditampilkan adalah rata-rata 5
pembacaan sehingga lebih representatif terhadap kondisi aktual beban listrik.

\subsection{Standar Deviasi sebagai Early Warning}

Nilai $\sigma$ yang meningkat sebelum arus bocor mencapai \emph{threshold}
memberikan potensi \emph{early warning} tambahan. Dalam skenario di mana arus
bocor meningkat bertahap (misalnya isolasi kabel yang mulai rusak), $\sigma$
akan meningkat lebih dulu sebelum $I_{\text{bocor}}$ melewati 15\,mA, memberi
waktu reaksi lebih awal bagi pengguna.

\subsection{Isolasi State Antar Cabang}

Salah satu pencapaian teknis penting dalam pengembangan ini adalah memastikan
bahwa injeksi kebocoran pada satu cabang tidak mempengaruhi \emph{state} cabang
lain. Hal ini dicapai dengan mengelola \texttt{current\_netral\_ma} secara
independen per sensor, dan hanya memodifikasi sensor yang ditunjuk tanpa
menyentuh data sensor lainnya.

\subsection{Keterbatasan Sistem}

\begin{itemize}
    \item \textbf{Simulasi bukan data \emph{hardware} nyata:} Nilai arus pada mode simulasi dihasilkan secara deterministik berdasarkan \emph{seed}, bukan dari pembacaan ADC ESP32 sesungguhnya. Implementasi \emph{hardware} nyata memerlukan integrasi driver ADC dan kalibrasi SCT-013.
    \item \textbf{Interval waktu tidak akurat:} Akumulasi energi mengasumsikan setiap pembacaan setara 1 jam, yang tidak mencerminkan waktu sesungguhnya antara dua pemanggilan \texttt{update\_all()}.
    \item \textbf{Tidak ada penyimpanan data:} Data monitoring hilang ketika program dihentikan. Implementasi lanjutan dapat menambahkan \emph{logging} ke file CSV atau \emph{database}.
\end{itemize}

% ══════════════════════════════════════════════════════════════════════════════
\section{Kesimpulan}
% ══════════════════════════════════════════════════════════════════════════════

\begin{itemize}
    \item Sistem ELISA berhasil diimplementasikan menggunakan bahasa pemrograman Rust dengan arsitektur modular berbasis OOP, mencakup \texttt{struct} Sensor, Controller, Actuator, dan MonitoringSystem beserta \emph{method}-method yang diperlukan.

    \item Prinsip deteksi kebocoran arus dengan formula $I_{\text{bocor}} = I_{\text{fasa}} - I_{\text{netral}}$ berhasil diimplementasikan dan terbukti bekerja sesuai ekspektasi pada seluruh skenario pengujian, dengan \emph{threshold} yang mengacu pada PUIL 2011.

    \item Empat metode komputasi numerik berhasil diterapkan: kalkulasi arus bocor (diferensial), Moving Average MA-5 (\emph{smoothing} sensor), standar deviasi sampel (monitoring kestabilan), dan akumulasi energi (integrasi Euler).

    \item \emph{Dashboard} web \emph{real-time} berhasil diintegrasikan menggunakan \texttt{tiny\_http} dan \emph{concurrency} \texttt{Arc<Mutex<MonitoringSystem>{}>}, memungkinkan pemantauan dari \emph{browser} tanpa menghentikan input terminal.

    \item Seluruh \emph{bug} kritis yang ditemukan selama pengembangan -- termasuk masalah isolasi \emph{state} antar cabang saat injeksi, kegagalan validasi pasca-reset simulasi, dan kondisi \emph{baseline} yang salah -- berhasil diidentifikasi dan diperbaiki secara sistematis.

    \item Sistem ini siap dikembangkan lebih lanjut ke arah integrasi \emph{hardware} ESP32 + SCT-013 untuk \emph{deployment} nyata di rumah tangga atau UMKM.
\end{itemize}

\end{document}
